\subsection{复射影线性簇的空间}

采用第六章 § 3 第 5 号中回忆过的记号，我们类似地定义复射影空间中射影线性簇的空间：

定义 2 。商空间 \(\mathbf{P}_{n,p}(\mathbf{C})\) ——拓扑空间 \(\mathbf{L}_{n+1, p+1}(\mathbf{C})\) 对等价关系 \(\Delta_{n,p}(\mathbf{C})\) 的商——称为 \(p \geqslant 0\) 维射影线性簇（在射影空间 \(\mathbf{P}_n(\mathbf{C})\) 中的）的空间。

由第六章 § 3 第 5 号命题 6 的论证，我们首先看出 \(\mathbf{P}_{n,p}(\mathbf{C})\) 是豪斯多夫的。其次，我们通过把（第六章 § 3 第 5 号命题 7 证明中的）子空间 \(\mathrm{V}_{n+1,p+1}\) 换成子空间 \(\mathrm{W}_{n+1,p+1}\) （ \(\mathrm{L}_{n+1,p+1}(\mathbf{C})\) 中由构成其所生成向量子空间的标准正交埃尔米特基的 \(p+1\) 个向量组所组成的）来证明它是紧的；这就是说，\(\mathrm{W}_{n+1,p+1}\) 由满足下列条件的矩阵 \(X=(x_{ij})\) 组成

\[
\sum_ {j = 0} ^ {n} x _ {i j} \bar {x} _ {i j} = \mathrm{I} \quad (\mathrm{I} \leqslant i \leqslant p + \mathrm{I}),
\]

\[
\sum_ {j = 0} ^ {n} x _ {i j} \overline {{{x}}} _ {k j} = 0 \quad (i \neq k).
\]

第六章 § 3 第 5 号命题 8 的证明对空间 \(\mathbf{P}_{n,p}(\mathbf{C})\) 原样适用，并表明这个空间是连通且局部连通的，且它的每一点都有一个同胚于 \(\mathbf{C}^{(p+1)(n-p)}\) 的邻域。最后，第六章 § 3 第 6 号命题 9 的证明也原样适用，因而格拉斯曼流形 \(\mathbf{G}_{n,p}(\mathbf{C})\) 同胚于 \(\mathbf{P}_{n,p}(\mathbf{C})\) 。

注。实与复的数空间（相应地，实与复的射影空间）共有的大部分性质，对以同样方式在四元数体 H 上定义的数空间（相应地射影空间）也成立；事实上，它们还可以推广到许多别的拓扑域和拓扑除环上去（参见习题 2 与 6）。

\ExerciseGroupsStart
% semantic-fragments: legacy-input-seam
\relax{}
